\section{How much of the proof does this replace?}

The conceptual saving is substantial and quite specific. The old tail
route first controls the maximizing ratio, then feeds that score interval
through moment bounds, charges the excess, and improves the score cap.
In the fixed-target route, the coefficient is $c$ from the start. The
nonnegative excess is retained on the left of the conservation identity.
One therefore seeks physical information about the same law directly.
The present result supplies an initial central gap by a short sequence of
scalar estimates instead of assuming it as an input to a refinement.

For the selected short presentation, use the signwise moment bound, the
fifth-moment square, the existing quadratic phase price, and one ordered
moment-circle update. The stronger full-phase argument in the appendix
is useful when a single pass is wanted, but need not be a compulsory step
in that presentation. This is a choice between proofs of the same
initial geometric conclusion; neither proof claims the later much
smaller defect and noise parameters.

The current main manuscript has 158 pages. It already uses the
least-counterexample reduction, so that earlier saving must not be
counted again. The sections devoted to complete Fourier estimates,
resource acquisition, confinement and count comparison occupy roughly
pages 33--89; finite support and the finite-index analysis occupy roughly
pages 101--156. These are measured heading spans, not disjoint estimates
of deletable pages. Most of their conclusions remain necessary.
Removing score interpolation alone saves a few local passages and a
layer of reasoning, not half of the whole proof.

A large further reduction needs a theorem that replaces a substantial
block, rather than a smaller numerical constant in each block. The most
promising target is a uniform estimate for the \emph{whole} fixed-target
remainder, with its actual-gap baseline and Fourier error, that reaches
the final moment/noise domain directly. A second is a strict count
comparison valid over that whole domain and down to arbitrarily small
positive noise. Those results could replace repeated acquisition or
separate local witnesses. They are still open here. The independent
finite base remains another essential task; moving its proof into a
separate paper would shorten one document but not the overall proof.

In a consolidated manuscript, replace the old score-dependent derivation
with the selected fixed-target derivation and prove each ingredient once.
Do not append all successive research notes as additional required
arguments. A readable organization is: finite base and first failure;
attainment/contact/support; geometry and conservation; one complete
Fourier engine; one ordered acquisition; count comparison; strict
exclusion. The new gap theorem advances the acquisition step. It does
not eliminate the Fourier engine, finite base or strict exclusion.
